\section{从障碍函数到原始---对偶内点法}
\label{sec:09-Convex-Optimization/06-Optimization-Algorithms/05}

一个资源分配问题摆在面前：三百个决策变量，八百条不等式约束，其中有预算上限、单项占比上限、几组两两之间的比例关系。把它交给 \texttt{cvxpy}，几百毫秒返回结果，日志显示只跑了二十几次迭代，输出里除了最优解还附带每条约束的对偶变量，可以直接读成放宽这条约束能省多少钱，见\hyperref[sec:09-Convex-Optimization/05-Duality/04]{``对偶变量是约束的影子价格''}。

上一篇的投影梯度法在这里用不上。它要求每步把点投回可行域，而投影到八百条不等式围成的多面体上，本身就是一个和原问题同等规模的二次规划——为了走一步，先要解一个一样难的问题。求解器显然没有走这条路。

它走的是另一条：不去把越界的点拉回来，而是让迭代从一开始就待在可行域内部，并且永远出不去。

\subsection{障碍把约束搬进目标}

标准形式写作

\[
\begin{aligned}
\text{minimize}\quad &f_0(x)\\
\text{subject to}\quad &f_i(x)\le0,\quad i=1,\ldots,m,\\
&Ax=b,
\end{aligned}
\]

其中 \(f_0,\ldots,f_m\) 凸、\(A\) 行满秩。不等式约束的麻烦在于它是硬的：可行与不可行之间没有过渡，梯度信息也不告诉你离哪条边界还有多远。

对数障碍把这道墙换成一个陡峭的斜坡。在严格可行点上定义

\[
\phi(x)=-\sum_{i=1}^m\log\bigl(-f_i(x)\bigr).
\]

每个约束贡献一项，当 \(f_i(x)\) 从负值升向 \(0\) 时该项趋于 \(+\infty\)。于是任何以最小化为目标的迭代都被自动挡在边界以内，不需要额外的可行性检查。代价是问题被改写了：把障碍按权重加进目标，得到

\[
\min_{Ax=b}\ \bigl\{t\,f_0(x)+\phi(x)\bigr\},\qquad t>0 .
\]

这是一个只带等式约束的光滑问题，Newton 法直接可用。参数 \(t\) 调节两者的相对重要性：\(t\) 小时障碍主导，解被推向可行域的``正中间''；\(t\) 大时原目标主导，解允许贴近边界。

\subsection{中心路径是一族被扰动的 KKT 解}

把 \(t\) 从小到大连续变化，每个 \(t\) 的解 \(x^\star(t)\) 描出一条曲线，称为中心路径。这条曲线不是随手取的插值，它每一点都满足一组可以写下来的方程。

对等式约束引入乘子 \(\tilde\nu\)，障碍子问题的驻点条件是

\[
t\nabla f_0(x)+\sum_{i=1}^m\frac{1}{-f_i(x)}\nabla f_i(x)+A^{\trans}\tilde\nu=0 .
\]

两边除以 \(t\)，并令

\[
\lambda_i=\frac{1}{-t\,f_i(x)}>0,\qquad \nu=\frac{\tilde\nu}{t},
\]

上式就变成 \(\nabla f_0(x)+\sum_i\lambda_i\nabla f_i(x)+A^{\trans}\nu=0\)，正是 KKT 的稳定性条件。原始可行、对偶可行也都自动成立。唯一的差别在互补松弛：由 \(\lambda_i\) 的定义直接得到

\[
\lambda_i\bigl(-f_i(x)\bigr)=\frac{1}{t},\qquad i=1,\ldots,m,
\]

而真正的 KKT 要求右端为 \(0\)，见\hyperref[sec:09-Convex-Optimization/05-Duality/03]{``KKT 把最优性写成一组方程''}。中心路径因此是一族被扰动的 KKT 解，扰动量恰好是 \(1/t\)，随 \(t\) 增大而消失。

\begin{center}
\begin{tikzpicture}[x=1cm,y=1cm,>=stealth]
  \fill[black!6] (0,0) -- (3.4,0) -- (4.2,1.6) -- (2.6,3.0) -- (0.4,2.4) -- cycle;
  \draw[black!55] (0,0) -- (3.4,0) -- (4.2,1.6) -- (2.6,3.0) -- (0.4,2.4) -- cycle;
  \node[font=\footnotesize,black!55] at (1.25,0.75) {可行域};
  \draw[black!80,thick] plot[smooth] coordinates
    {(2.10,1.50) (2.55,1.28) (2.90,1.00) (3.15,0.70) (3.32,0.40) (3.39,0.12)};
  \fill[black!80] (2.10,1.50) circle (1.7pt);
  \fill[black!80] (2.90,1.00) circle (1.7pt);
  \fill[black!80] (3.32,0.40) circle (1.7pt);
  \fill[black] (3.4,0) circle (2.2pt);
  \node[above left,font=\footnotesize,black!75] at (2.18,1.50) {解析中心};
  \node[below right,font=\footnotesize] at (3.44,0.04) {\(x^\star\)};
  \node[font=\footnotesize,black!70] at (6.0,1.10) {\(t\) 增大，中心路径滑向边界};
  \draw[->,black!55] (4.62,1.00) -- (3.46,0.52);
  \draw[->,black!60] (0.30,3.55) -- (1.02,2.95);
  \node[left,font=\footnotesize,black!60] at (0.24,3.55) {\(-c\)};
\end{tikzpicture}

\emph{障碍在边界处发散，把迭代锁在可行域内部；\(t\) 从小到大，解从解析中心出发，沿一条光滑曲线走向边界上的最优点。}
\end{center}

图里那条曲线的两端各由一项主导。\(t\to0\) 时目标几乎只剩 \(\phi\)，极小点是使各条约束的松弛量乘积最大的那个点，也就是离所有边界都``尽量远''的解析中心；它与目标函数无关，只由可行域的形状决定。\(t\) 增大后原目标开始拽动这个点，路径朝 \(-c\) 的方向滑，最终顶到边界上的最优顶点。关键在于中间的每一点都是内点——路径从不触碰边界，只是无限逼近。这也说明为什么求解器输出的解通常不会精确落在顶点上，而是落在离顶点极近的内部，需要时再做一次识别活跃约束的后处理。

\subsection{外层的账可以提前算清}

沿中心路径构造出的 \(\lambda,\nu\) 是一组合法的对偶可行解，把它代入 Lagrangian 对偶函数，就得到一个真实的下界，见\hyperref[sec:09-Convex-Optimization/05-Duality/01]{``Lagrangian 怎样制造全局下界''}。算一下这个下界与当前目标值的距离：

\[
f_0(x)-g(\lambda,\nu)=-\sum_{i=1}^m\lambda_if_i(x)=\frac{m}{t}.
\]

对偶间隙是 \(m/t\)，一个不含任何未知量的表达式。要把次优度压到 \(\varepsilon\) 以下，只需

\[
t\ge\frac{m}{\varepsilon}.
\]

这是内点法最实用的性质之一：还没开始算，就知道要把 \(t\) 推到多大；算完之后交出的不只是一个解，还有一份可以被独立核验的精度证书，这正是\hyperref[sec:09-Convex-Optimization/05-Duality/02]{``对偶间隙何时闭合''}所说的那种交付物。

于是算法结构就定了。外层每次把 \(t\) 乘以一个大于 \(1\) 的因子（常取 \(10\) 到 \(100\)），内层用 Newton 法把当前子问题解到中心路径附近，上一个中心点作为下一个子问题的初值——这叫路径跟踪。外层的轮数是 \(\log(m/(\varepsilon t_0))\) 除以 \(\log\mu\)，对精度只是对数依赖，把 \(\varepsilon\) 再缩小一千倍也只多跑几轮。

内层的轮数则由自协调性兜底。对数障碍属于自协调函数：它的三阶导数被自身的 Hessian 按一种与坐标系无关的方式控住，于是上一篇\hyperref[sec:09-Convex-Optimization/06-Optimization-Algorithms/02]{``Newton 法用曲率校正方向''}里那个 Newton decrement 不只是停止判据，还能给出从当前点到子问题最优值的可计算上界，进而给出所需 Newton 步数的显式界——这个界不依赖问题的条件数，也不依赖变量的量纲，正因为 Newton 步和 decrement 都是仿射不变的。

两层合起来，总迭代次数是问题规模和 \(\log(1/\varepsilon)\) 的多项式。这一点的分量值得说清楚：在此之前，``凸''只是一个保证局部即全局的结构性质，并不自动意味着可解；内点法把它兑换成了一般凸问题的多项式时间求解框架，也把线性规划、二次规划、二阶锥规划、半正定规划统一在同一套机器下面，见\hyperref[sec:09-Convex-Optimization/04-Convex-Problems/01]{``标准形式与线性规划''}与\hyperref[sec:09-Convex-Optimization/04-Convex-Problems/04]{``半正定规划优化矩阵不等式''}。凸与非凸之间那条计算上的分界线，很大程度上是被这个框架划出来的，见\hyperref[sec:09-Convex-Optimization/09-Complexity-and-Approximation/01]{``P、NP 与 NP-hard：给问题按计算难度分类''}。

\subsection{原始---对偶方法直接追踪扰动后的 KKT}

障碍方法把对偶变量当作副产品，每次内层解完才反推出来。现代求解器更进一步，把 \(x,\lambda\) 和松弛变量放在一起同时更新。引入 \(s_i>0\) 把不等式写成等式 \(f_i(x)+s_i=0\)，KKT 的互补条件是 \(\lambda_is_i=0\)，内点法把它松弛为

\[
\lambda_is_i=\mu,\qquad \mu>0,
\]

再对由原始残差、对偶残差和这组互补残差拼成的非线性方程组做 Newton 迭代，同时让 \(\mu\) 逐步下降。\(\mu\to0\) 时轨迹收敛到真实的 KKT 点。

这种写法有几处实际好处。原始与对偶的可行性残差都摆在眼前，随时可以判断是收敛慢还是问题本身不可行；\(\mu\) 不必与当前的互补间隙严格绑定，predictor--corrector 之类的策略可以先试一个激进的方向，再用一次校正把偏离中心路径的部分补回来，实测能把迭代数再压掉三成；最重要的是每步要解的那个 KKT 线性系统保持了原问题的稀疏模式，可以做一次符号分析、多次数值分解，成本远低于按稠密矩阵处理。这一步仍然是解方程而不是求逆，用的是对称不定矩阵的分解，见\hyperref[sec:08-Numerical-Analysis/05-Numerical-Linear-Algebra/03]{``Cholesky：利用正定结构''}。

\subsection{障碍不是惩罚}

容易与障碍混淆的是外部惩罚法。惩罚项形如 \(\rho\sum_i\max\set{0,f_i(x)}\)，它只在约束被违反时才产生代价，允许迭代暂时不可行，靠增大 \(\rho\) 把违规压回去。障碍恰好相反：它只在严格可行域内部有定义，出界就是 \(+\infty\)。两者的迭代轨迹一个在外一个在内，对初值的要求也不同——惩罚法可以从任意点起步，障碍法必须从一个严格可行点起步。

这个起步点未必现成。找不到时要先解一个可行性辅助问题：引入一个标量 \(s\)，最小化 \(s\) 使得 \(f_i(x)\le s\) 对所有 \(i\) 成立，这个问题自带严格可行起点。若最优的 \(s<0\)，就得到了严格可行点；若最优的 \(s>0\)，则原问题不可行，而且这个结论同样带着对偶证书，不是``算不出来''而是``证明了不存在''。

\subsection{它适合什么规模的问题}

内点法的迭代数少而每步昂贵，这个组合决定了它的适用面。几十次迭代就能把间隙压到 \(10^{-8}\)，条件数再差也不太影响轮数；但每一轮都要分解一个规模与约束数相当的线性系统，稠密时是 \(O(n^3)\)，稀疏时取决于填充量。中等规模、结构化、要求高精度并需要对偶证书的问题最适合它——投资组合、最优控制的求解子问题、SVM 的对偶形式都属此列，见\hyperref[sec:09-Convex-Optimization/07-Applications/04]{``风险、收益与投资组合''}。变量上千万时分解本身就装不下内存，那时该退回一阶方法或做问题分解，用精度换规模。

所以``迭代次数少''不能直接读成``总时间短''。这一单元从梯度下降一路走到内点法，五种方法的差别归根到底不在名字上：能拿到的是梯度、次梯度、近端算子还是完整的 Hessian 与稀疏结构，愿意为一步付出多少，需要多少位有效数字，这几件事一定，方法基本就定了。

算法讲到这里，剩下的问题是它们落在真实问题上会长成什么样。\hyperref[ch:077]{``建模案例：凸性怎样改变真实问题''}把经验风险、\(L_1\) 稀疏、最大间隔与投资组合逐一拆开，看凸性在建模阶段究竟买到了什么。
